%\chapter{ABSTRACT}
\clearpage
\phantomsection
\addcontentsline{toc}{chapter}{ABSTRACT}
\begin{center}
\textbf{\large  ABSTRACT}\\[0.5cm]
\end{center}

This thesis presents a comprehensive investigation into the aerothermal characteristics of Hypersonic Inflatable Aerodynamic Decelerators (HIADs), with a particular focus on the impact of their scalloped surface topology on flow physics and aerodynamic heating. Utilizing the Direct Simulation Monte Carlo (DSMC) method implemented in the SPARTA solver, the study develops a computational framework for rarefied gas dynamics analysis of HIAD reentry, where the freestream is near-continuum ($Kn \approx 3.4\times10^{-5}$ on a 3\,m body at 52\,km) and the shock layer locally spans the transitional and near-continuum regimes ($10^{-5} \lesssim Kn \lesssim 1.2\times10^{-2}$). The framework integrates parametric CadQuery geometry generation, DSMC aerothermal simulation, and simple Genetic Algorithm optimization driven by a closed-form cost. The optimization runs in two passes: a DSMC re-simulation of the first winner sets a bounded correction factor for the analytic heat load, and the genetic algorithm is re-run with the corrected cost before a second, final DSMC validation. The first-pass winner fails the heat gate at 175.580\,J/cm\textsuperscript{2} against a 165.716\,J/cm\textsuperscript{2} target, while the second-pass vehicle meets the same numeric target at its own trajectory state at 117.526\,J/cm\textsuperscript{2}. The methodology is grounded in the Boltzmann transport equation and checked against IRVE-3 flight data, with the Sutton-Graves stagnation-point correlation and Fay--Riddell theory serving as thermal benchmarks; the single-point Sutton--Graves gate estimate under-predicts the flight-integrated heat load by 15.0\% while the DSMC drag solution over-predicts the peak deceleration by 15.3\%, which are stated as limitations rather than absorbed into the validation claim. Key findings include the characterization of scalloped toroidal heating augmentation, Knudsen number field mapping across transitional and continuum regimes, and the demonstration of a four-level grid independence study (70$\times$70, 100$\times$100, 139$\times$139 and 200$\times$200 cells) supporting the $139\times139$ production grid, which is retained on cost grounds: a five-level study ($70\times70$, $100\times100$, $139\times139$, $200\times200$ and $250\times250$ cells) leaves neither the drag nor the heat-flux series grid-converged, and the peak surface heat flux carries a roughly $15\,\%$ grid band with a resolution-dependent location. The 5-species air model (N$_2$, O$_2$, NO, N, O) with Variable Soft Sphere collision physics captures real-gas effects including vibrational excitation and dissociation, enabling accurate prediction of aerothermal loads for HIAD survivability optimization.

\noindent \textbf{Keywords}: Hypersonic Inflatable Aerodynamic Decelerator (HIAD), Aerothermodynamics, Direct Simulation Monte Carlo (DSMC), SPARTA, Rarefied Gas Dynamics, Scalloped Geometry, Aerodynamic Heating, Genetic Algorithm Optimization